% sec-09: risks, stop conditions, and the cost of the conversion.  No figure.
% The cons of the chosen approach are stated here in full, not softened.
\section{Risks, Stop Conditions, and What This Is Not}

\subsection{What the Conversion Costs}

\draftinstr{The section a reviewer will test the document's honesty against.
Four numbered concessions, each in its own short paragraph, each stated without
mitigation in the first sentence and with the plan's answer in the second.
(1) The individual-scientist framing that made all ten prior applications
distinctive is abandoned, and it is the framing the report actually rewards.
(2) Capitalization and cost detail the author might prefer to keep private is
published here, permanently and under a DOI. (3) The Pioneer, HHMI and FRO
recipients are person-based or organization-based and will read a company plan
as off-mechanism, which is three of the ten prior recipients. (4) A financing
paper is not a scientific publication and fits no journal, which limits it to
Zenodo. Source: \appfile{funding/pdac-funding-applications/final-apply/sections/sec-02-ten-applications.tex}
for the ten recipients and their mechanism types.}

\subsection{The Risk Register}

\draftinstr{Set Table 18: eight to ten rows, five columns (risk, the milestone
or figure it attaches to, likelihood stated as low, moderate or high, the
consequence, and the mitigation already in place as opposed to planned). The
mitigation column must distinguish the two: a mitigation that exists today is
written plainly, and a mitigation that does not yet exist is written as "none
today" with the instrument that would create it. Rows must include at minimum:
the site agreement not executing; a clinical hold; two of three DLTs at dose
level 1; the Phase II award not being made; the founder falling below 50 percent
ownership; an investigator acquiring a \S54.2 interest; the advisory model being
mistaken for a directive; and the 81.9 credibility score failing to reproduce.
Sources: \appfile{funding/capitalization-plan/graphviz/fig-14-stop-condition-fault-tree.gv.md}
and \appfile{funding/pdac-funding-applications/final-apply/sections/sec-10-risks-and-limits.tex}.}

\subsection{What This Document Is Not}

\draftinstr{A short, plain list in prose, each item one sentence. It is not a
submission of record. It is not an agreement with UC San Diego, Moores Cancer
Center, or any other institution, and no agreement of any kind exists. It is not
a completed raise: no term sheet, SAFE or subscription agreement exists, and
every capitalization figure is a plan. It is not a claim that daraxonrasib is
first in human; the agent is investigational and already in Phase 3 evaluation.
It is not medical or regulatory advice. No patient has been treated. No robotic
configuration has been specified or cleared. Source:
\appfile{funding/potential-partners/UC-San-Diego/} for the positioning
correction the parent work's audit pass introduced.}

\subsection{The Two Single Points of Failure}

\draftinstr{Close the section by naming them, because Figure 14 makes them
visible and a reader who has skipped the figure should still meet them. H1, the
month-nine gate, fails on either a bench latency above 250 ms at the 95th
percentile or an IND clinical hold. H5, the ownership test, fails on either the
founder falling below 50 percent or an investigator acquiring a \S54.2 interest.
Neither is technical in the sense that more engineering fixes it, and both are
knowable in advance, which is the only reason they are acceptable.}
